% titel: Blatt 2, A9: Reaktionsgleichungen ausgleichen
% skript: Mortimer Kap. 4.6 (Reaktionsgleichungen, stoichiometrie)

% Teil I: Werkzeuge, kein Lösungsweg.

Eine Reaktionsgleichung ist dann korrekt, wenn die Anzahl der Atome jeder
Sorte auf beiden Seiten gleich ist (Erhaltung der Atome). Koeffizienten
(ganze Zahlen vor den Formeln) werden so gewaehlt, dass diese Bedingung
minimal erfüllt ist.

\begin{werkzeugkasten}[Ausgleichen einer Reaktionsgleichung]
Vorgehen:
\begin{enumerate}
  \item \textbf{Bilanzieren nach Elementsorten:} Für jede Atomsorte
        Zähler auf der Edukt- und der Produktseite gleichsetzen.
  \item \textbf{Nacheinander ausgeraeumt:} Zunaechst die Atomsorte, die nur in
        einer Verbindung je Seite vorkommt, dann die übrigen; Elemente, die
        als reine Elemente ($\mathrm{O_2}$, $\mathrm{C}$) auftreten, kommen
        zuletzt dran.
  \item \textbf{Koeffizienten ganzzahlig:} Mit einem geeigneten gemeinsamen
        Faktor multiplizieren, falls Bruchkoeffizienten auftreten; Ergebniswerte
        vielfach ganzen Zahlen.
\end{enumerate}
\textbf{Tabellenmethode} (nur bei Bedarf): Man setzt Buchstaben als unbekannte
Koeffizienten, stellt für jede Atomsorte eine Gleichung auf und löst.
Bei \textbf{Redox-Bruttoreaktionen} prüfen, ob die Gesamtladung (falls Ionen
beteiligt) ebenfalls erhalten bleibt.
\textit{Skript:} Mortimer Kap. 4.6 (Reaktionsgleichungen), Übung 2/9.
\end{werkzeugkasten}

\begin{achtung}[Darauf achten]
\begin{itemize}
  \item Korrekt schreiben: die Stoffe selbst dürfen \emph{nicht} veraendert
        werden --- nie die Summenformel, nur die Koeffizienten davor.
  \item Koeffizienten sind die kleinste ganze Zahlenmenge (ggT auf $1$).
  \item Atome des gleichen Elements, die in mehreren Verbindungen stecken,
        als Summe bilanzieren (z.~B. Wasserstoff in mehreren Produkten).
  \item Bei einer Verbindung, die als Gegenstand der Reaktion nur auf einer
        Seite auftritt, ist deren Koeffizient frei waehlbar (Startpunkt).
\end{itemize}
\end{achtung}
